\vspace{-0.2cm}
\section{Background}
\vspace{-0.2cm}

\textbf{Language Generation as an MDP.}
Language generation tasks can be modeled as a Markov Decision Process (MDP) defined by a tuple \((\mathcal{S}, \mathcal{A}, \mathbb{P}, \mathcal R)\), where \(\mathcal{S}\) is the state space, \(\mathcal{A}\) is the action space, \(\mathbb{P}\) represents transition dynamics, and \(R\) is the reward function. Specifically, at each time step \(t\), the state \(s_t \in \mathcal{S}\) consists of the prompt tokens \(\boldsymbol x\) along with previously generated tokens \(\boldsymbol y_{<t}=[y_1,\dots,y_{t-1}] \), that is, \( s_t = [\boldsymbol x, \boldsymbol y_{<t}]\). The action \(a_t \in \mathcal{A}\) corresponds to selecting the next token \(y_t\). The decision-making process follows a policy \(\pi_\theta(a_t|s_t)\), parameterized by \(\theta\), which defines the probability of the next token conditioned on the current state. The transition dynamics \(\mathbb{P}\) are deterministic: given state \( s_t \) and action (token) \( a_t \), the next state \( s_{t+1} \) is obtained by concatenating the selected token to the current state: \( s_{t+1} = \mathbb{P}(s_t, a_t) = [s_t, a_t] \). The reward function assigns 0 intermediate rewards, providing only a sparse binary reward \( \mathcal R(\boldsymbol x,\boldsymbol y) = 1 \text{ or } 0\) at episode termination, where $\boldsymbol y=[y_1, y_2, \dots, y_T]$ is the full generated response, indicating whether the generated response matches the ground truth or not. The value function \( V(s) \) represents the expected cumulative reward from state \( s \). The state-action value function \( Q(s,a) \) corresponds to the expected cumulative reward from choosing action \( a \) at state \( s \). The advantage function \( A(s,a) \), defined as \( Q(s,a)-V(s) \), measures the improvement in expected reward realized by taking an action \( a \) at state \( s \). Under deterministic transition dynamics and sparse binary rewards, the advantage function simplifies to \( A(s,a)=V(s')-V(s)\), where \( s' = \mathbb{P}(s,a) \) is the next state reached deterministically from state \( s \) after taking action \( a \).

\textbf{RL for LLMs.} PPO \cite{schulman2017proximalpolicyoptimizationalgorithms, ouyang2022traininglanguagemodelsfollow} and GRPO \cite{shao2024deepseekmathpushinglimitsmathematical} are two widely adopted reinforcement learning algorithms for optimizing large language models. PPO introduces a clipped surrogate objective to stabilize training by constraining policy updates near the previous policy:
\begin{equation}
\mathcal{J}_{\text{PPO}}(\theta) = \mathbb{E}_{\boldsymbol x \sim \mathcal{D}, \boldsymbol y \sim \pi_{\theta_{\text{old}}}(\cdot|\boldsymbol x)} \bigg[ \frac{1}{|\boldsymbol y|} \sum_{t=1}^{|\boldsymbol y|} \min\Big( r_t(\theta) \hat A_t, \mathrm{clip}(r_t(\theta), 1 - \epsilon, 1 + \epsilon) \hat A_t \Big) \bigg],
\label{eq:PPO}
\end{equation}
where \( r_t(\theta) \) is defined as $\frac{\pi_{\theta}(a_t|s_t)}{\pi_{\theta_{\text{old}}}(a_t|s_t)}$, and \(\epsilon\) is a small hyperparameter that limits excessively large updates. PPO estimates token-level advantages \(\hat A_t\) via GAE \cite{schulman2018highdimensionalcontinuouscontrolusing}, requiring a critic model to predict token-level values. Compared to PPO, GRPO eliminates the need for the critic model and instead estimates the trajectory-level advantages by normalizing each response's reward within the sampled group. These trajectory-level advantages are then assigned uniformly to all tokens in the corresponding trajectory to obtain \(\hat A_t\). Similar to PPO, GRPO adopts a clipped objective, together with a directly imposed KL penalty term. Recent work \cite{kazemnejad2024vineppounlockingrlpotential} highlights that accurately training the critic in PPO is challenging, and proposes VinePPO, a critic-free alternative that partitions the reasoning trajectory into discrete steps using heuristic rules (e.g., line breaks) and estimates step-level advantages through Monte Carlo (MC).

Compared to VinePPO, our proposed segment-level framework, SPO, adopts a more general concept of segmentation, allowing arbitrary partitions without enforcing semantic coherence. This enables us to freely adjust the granularity anywhere between token-level and trajectory-level, and also facilitates adaptation to broader and less structured tasks such as code generation. Moreover, we introduce a novel tree-based sampling method, where segment advantages are estimated concurrently with trajectory generation. This obviates the need for VinePPO’s costly resampling to estimate advantages, thereby substantially improving efficiency and making our approach effective in long CoT scenarios. Notably, VinePPO can be viewed as a special case within our more general SPO framework. More related work is provided in Appendix \ref{appendix:more_related_work}.